\chapter{Kesimpulan dan Saran}

\section{Kesimpulan}

Penelitian ini berhasil membangun \textit{pipeline} berbasis Python 
yang mengintegrasikan konversi otomatis kode PHP lama, pemindaian 
keamanan, dan pemetaan kepatuhan ISO/IEC 27001:2022 dalam satu alur 
kerja yang dapat dijalankan sepenuhnya secara lokal. Berdasarkan 
pengujian pada dua studi kasus nyata, berikut kesimpulan yang 
menjawab tujuan penelitian.

\textbf{Pertama}, konversi otomatis PHP berbasis Rector berhasil 
diimplementasikan dengan tingkat keberhasilan 97,6\% pada dataset 
FT UGM dan 91,7\% pada dataset CBT DTI UGM, dengan validitas 
sintaks pasca-konversi mencapai 100\% pada keduanya. Kegagalan 
pada dua \textit{file} di kedua dataset disebabkan oleh 
\textit{bug} internal Rector, bukan dari kode sumber.

Validasi berbasis \textit{sampling} terhadap 12 \textit{file}
aplikasi menunjukkan rata-rata \textit{similarity score} 0,9214
antara keluaran Rector dan referensi konversi manual,
mengindikasikan tingkat kesesuaian yang tinggi antara
pendekatan otomatis dan manual.

\textbf{Kedua}, mekanisme pemeriksaan keamanan yang dipetakan ke
kontrol ISO/IEC 27001:2022 Annex A berhasil diimplementasikan secara
deterministik dan dapat diaudit sepenuhnya. Nilai $\Delta F = 0$ pada
kedua studi kasus mengonfirmasi bahwa konversi Rector tidak
memperkenalkan kerentanan baru ke dalam basis kode, sebuah konfirmasi
yang penting mengingat transformasi otomatis berskala besar berpotensi
mengubah perilaku kode secara tidak terduga. Status
\texttt{NON\_COMPLIANT} yang dihasilkan bukan cerminan kegagalan
\textit{pipeline}, melainkan potret kondisi nyata kode warisan yang
memang sudah mengandung permasalahan keamanan sebelum dilakukan
remediasi. Seluruh hasil analisis ini disatukan dalam laporan 
JSON yang dibuat secara otomatis pada setiap eksekusi \textit{pipeline}. 
Laporan otomatis ini dapat dimanfaatkan langsung sebagai artefak teknis pendukung dalam proses audit sertifikasi standar keamanan tanpa memerlukan rekapitulasi manual tambahan.

\textbf{Ketiga}, komponen validasi PHPStan memberikan manfaat konkret
dengan mengekspos ribuan potensi masalah kompatibilitas tipe yang
tidak akan terdeteksi sebelum kode benar-benar dijalankan di
lingkungan produksi. Mayoritas temuan merupakan \textit{technical
debt} yang sudah ada jauh sebelum migrasi dilakukan dan baru terlihat
karena sistem tipe PHP 8.3 jauh lebih ketat dibanding versi
sebelumnya. Hal ini menunjukkan bahwa \textit{pipeline} 
tidak hanya menjalankan transformasi sintaksis, tetapi juga 
memetakan kondisi aktual basis kode yang sebelumnya tidak 
terdeteksi akibat sifat permisif sistem tipe pada PHP versi lama.

\textbf{Keempat}, \textit{AI Engine} berbasis LLM \textit{open source}
lokal terbukti menghasilkan keluaran yang secara terukur lebih
informatif dibandingkan keluaran \textit{pipeline} tanpa komponen AI.
Eksperimen \textit{blind rating} oleh dua rater independen dari DTETI
FT UGM terhadap 11 temuan Semgrep menunjukkan bahwa kondisi T1
(dengan \textit{AI Engine} aktif) unggul atas kondisi B0 (tanpa
\textit{AI Engine}) pada dimensi Kejelasan (+0,591 poin) dan
\textit{Actionability} (+0,500 poin) dari skala Likert 7. Perbedaan
rata-rata keseluruhan sebesar 0,284 poin dikonfirmasi signifikan
secara statistik melalui uji Wilcoxon \textit{signed-rank} satu arah
($W = 44{,}0$, $p = 0{,}047$, $n = 10$ temuan SSRF). Sebaliknya,
tidak ada perbedaan pada dimensi Akurasi maupun Ketepatan pemetaan
ISO, yang menunjukkan bahwa pemetaan ISO berbasis aturan deterministik
yang telah diimplementasikan sudah memadai dan tidak memerlukan
penguatan dari model bahasa. Hasil ini perlu dibaca dengan
kehati-hatian mengingat nilai Cohen's kappa antar-rater yang rendah
($\kappa = {-0{,}097}$), yang mengindikasikan perbedaan gaya respons
antara kedua rater.

Sebagai bagian dari evaluasi ini, komparasi performa lima model LLM
\textit{open source} menunjukkan bahwa \textit{Format Compliance Rate}
dan \textit{Syntax Validity Rate} merupakan dua dimensi yang independen
dan tidak berkorelasi. Qwen2.5-Coder direkomendasikan sebagai model
utama karena secara konsisten menghasilkan SVR 100\% di seluruh
kondisi eksperimen. SVR terbukti sebagai metrik evaluasi yang paling
sahih untuk sistem berbasis LLM lokal karena hasilnya dapat
diverifikasi secara independen menggunakan \texttt{php -l}, berbeda
dari \textit{confidence score self-reported} yang tidak berkorelasi
dengan kualitas kode yang sebenarnya dihasilkan.


\section{Saran}

Berdasarkan keterbatasan yang teridentifikasi, berikut arah 
pengembangan yang disarankan untuk penelitian selanjutnya.

\begin{enumerate}
    \item Evaluasi LLM perlu diperkuat dari dua sisi: 
    validasi semantik menggunakan \textit{test suite} otomatis 
    untuk memverifikasi kebenaran logika perbaikan, dan 
    perluasan dataset dengan keragaman jenis kerentanan yang 
    lebih tinggi agar kesimpulan perbandingan model dapat 
    digeneralisasi lebih luas. Sebagai kelanjutan jangka 
    panjang, eksplorasi \textit{fine-tuning} dengan LoRA atau 
    QLoRA dapat dilakukan dengan menjadikan penelitian ini sebagai \textit{baseline}.

    \item Akurasi laporan PHPStan pada basis kode
CodeIgniter 3.x dapat ditingkatkan melalui dua pendekatan:
(1) \textit{custom stubs} yang mendefinisikan tipe untuk
\textit{magic property} CI3 seperti \texttt{\$this->db} dan
\texttt{\$this->load} secara eksplisit, sehingga PHPStan
tidak perlu bergantung pada \textit{runtime resolution};
dan (2) penggantian mekanisme \textit{noise filter} berbasis
\textit{substring matching} yang digunakan saat ini dengan
konfigurasi \textit{baseline} PHPStan yang lebih akurat dan
tidak berisiko menyaring temuan nyata secara tidak sengaja.


    \item Validasi pemetaan ISO/IEC 27001:2022 oleh 
    auditor atau praktisi bersertifikat perlu dilakukan untuk 
    meningkatkan kredibilitas laporan kepatuhan yang dihasilkan.

    \item Integrasi ke CI/CD dan perluasan ke bahasa 
    pemrograman lain dapat dieksplorasi lebih lanjut. Arsitektur modular 
    \textit{pipeline} memudahkan penggantian komponen Rector dan 
    Semgrep dengan alat setara untuk Python, JavaScript, atau 
    bahasa lain, sekaligus membuka kemungkinan integrasi ke 
    GitHub Actions atau GitLab CI.

    \item Pengujian terhadap PHP 8.4 dan 8.5 serta
    perbandingan dengan alat komersial seperti SonarQube atau
    Snyk pada dataset yang sama perlu dilakukan untuk memvalidasi
    cakupan dan akurasi \textit{pipeline} secara lebih komprehensif.

    \item Validasi efektivitas \textit{AI Engine} perlu direplikasi
    dengan jumlah rater yang lebih banyak dan desain eksperimen yang
    mengendalikan \textit{position bias}. Nilai Cohen's kappa yang
    rendah ($\kappa = {-0{,}097}$) pada eksperimen \textit{blind
    rating} mengindikasikan perbedaan gaya respons antar-rater yang
    substansial, sehingga kesimpulan tentang keunggulan T1 atas B0
    memerlukan konfirmasi lebih lanjut. Selain itu, keberagaman dataset
    yang lebih tinggi -- mencakup lebih banyak jenis kerentanan selain
    SSRF -- diperlukan agar hasil validasi dapat digeneralisasi ke
    konteks yang lebih luas.
\end{enumerate}